\section{1. Introduction}\label{introduction}

\textbf{Research question.} How does Federal Reserve communication move asset prices, and when does it move the portfolios of households and institutions? We answer with a decomposition of monetary-policy news into a conventional component (MP) and an information component (CBI) --- constructed by letting financial markets invert the central bank's statement under sign restrictions --- and by tracing each component into (i) announcement-day prices and (ii) monthly mutual-fund flows.

\textbf{The theoretical anchor (Pillar 1, rational expectations).} CBI is not a property of the text; it is the market's Bayesian update about the central bank's private information. The sign-restriction identification \emph{is} the inversion rule: MP shocks move rates and equity premia with the classical signs; CBI shocks (Coibion-Gorodnichenko-Ulig style information surprises) move them with correlated but distinct signs. {[}TODO: 1-2 paragraphs on why \emph{statements} --- not just decisions --- are the object; JK (2020) precedent.{]}

\textbf{What we find.} 1. \emph{Prices} react to both components, but the dictionary-measured \emph{tone} of the statement matters for prices only in the forward-guidance era (beta = −1.01, p = 0.005, R² = 32.8\% in the FG subsample; near-zero outside it). Tone's price effect is concentrated exactly where communication paradigm shifted. 2. \emph{Portfolios} do not: neither shocks nor tone nor LLM-scored text add explanatory power to monthly fund flows (H4: p = 0.64 raw; p = 0.17 LLM). This null is the paper's ruling result --- ``prices first, portfolios later.'' 3. The risk-taking channel runs through \emph{good} news, not bad: positive CBI shocks (central bank reveals better information) drive de-risking --- flows out of high-risk classes when Risk × CBI interact (Risk×CBI beta = −0.066, p = 0.036; government-bond CBI beta = +0.76, p = 0.037, same-month window). This reverses the standard ``bad news → reach for safety'' reading: it is \emph{good news about the economy} that triggers portfolio rebalancing toward safety. {[}TODO: connect to APT/Pillar 3 --- risk-premium repricing, not risk appetite.{]} 4. \emph{Measurement}: the LM dictionary fails by compression (not by mis-reading) exactly when the Fed's communication paradigm shifts --- on the 11 divergence statements the dictionary output is nearly constant (std 1.10) while multi-model LLM scoring separates directions sharply (std ≈ 70); sign agreement across four vendor families is 100\% on those statements. We position the LLM measure as a diagnostic of dictionary failure, not a better scorer (it adds no price or flow predictiveness).

\textbf{Contribution.} {[}TODO: 3 bullets --- (i) inversion-rule identification for statements; (ii) the flow null as evidence for price-first transmission with a \emph{theoretical} reason (timing structure: flows are monthly, prices are instantaneous); (iii) measurement-failure evidence that disciplines the tone literature.{]}

\textbf{Roadmap.} {[}TODO: 1 paragraph.{]}

\section{2. Theoretical Framework}\label{theoretical-framework}

\subsection{2.1 Pillar 1 --- Rational expectations and the inversion rule}\label{pillar-1-rational-expectations-and-the-inversion-rule}

The central bank reveals private information through statements; markets invert it. Formally, the announcement is a signal in a Kanoidal/CGU structure; the identified CBI shock is the market's posterior revision. Sign restrictions are not an assumption bolted on --- they are the observable implication of the model. {[}TODO: formal model section --- 2-period signal extraction, signs table.{]}

\subsection{2.2 Pillar 2 --- Information economics of policy news}\label{pillar-2-information-economics-of-policy-news}

Kuttner (2001) surprise construction; Bauer-Swanson (2022) rejection of high-frequency calibration short-cuts; Jarociński-Karadi (2020) MP/CBI decomposition. We extend the decomposition to \emph{statement text} (CBI from statements, 60 identified meetings) alongside decision-window shocks (91 MP meetings). {[}TODO: expand with the exact identification equations; codebook anchor Table 2.{]}

\subsection{2.3 Pillar 3 --- APT: why portfolios should lag prices}\label{pillar-3-apt-why-portfolios-should-lag-prices}

Under arbitrage pricing, professional arbitrageurs absorb announcement-day price moves instantly; household/institutional portfolios adjust through slower channels (rebalancing rules, flows, lags). The testable implication is the \emph{timing structure}: price effects at zero horizon, flow effects (if any) at monthly horizon, and any tone/text effect on flows must vanish conditional on the shocks. {[}TODO: formalize as 3-type-agent model; state the null precisely.{]}

\textbf{De-risking prediction.} Good CBI news lowers the risk premium → marginal investor's required compensation falls → high-risk assets are \emph{over-valued at the margin} → rule-based portfolios rotate toward safety. Prediction: Risk × CBI(+) \textless{} 0 for risky-class flows, CBI(+) \textgreater{} 0 for government-bond flows. {[}TODO: derive formally.{]}

\section{3. Data and Measurement}\label{data-and-measurement}

\subsection{3.1 Statements and shocks}\label{statements-and-shocks}

161 FOMC statements, 2006--2026 (scrubbed: \texttt{data/fomc\_statements\_clean.json}). Decision-window shocks (91 MP / 60 CBI identified meetings) from fed-funds futures + 2y/10y windows {[}TODO: exact windows{]}. CBI-from-statements via sign restrictions {[}TODO: spec{]}.

\subsection{3.2 Fund flows}\label{fund-flows}

Morningstar-class monthly net flows (7 classes: government bonds, corporate bonds, real assets, large-cap/DM/EM/small-cap equity), N = 157 per class-month regression. Flow gap rule ≤ 3 calendar months (server-side aggregation, point-in-time classification, pseudo B-S removed). {[}TODO: data providers, cleaning log.{]}

\subsection{3.3 Text measures}\label{text-measures}

\begin{itemize}
\tightlist
\item
  LM tone: per-100-words hawkish-dovish; never called ``sentiment'' (codebook discipline).
\item
  LLM score: full-text stance read by multi-vendor committee (GPT-4.1-mini, Claude-haiku-4.5, DeepSeek-flash, Gemini-2.5-flash; majority sign rule), rebuilt from statement text --- \emph{not} meeting metadata (the earlier proxy was circular; see §7).
\item
  Variables: one name, one definition (codebook.docx); ``shock'' reserved for innovations.
\end{itemize}

\section{4. Prices: The FG Channel}\label{prices-the-fg-channel}

\textbf{Headline.} Tone predicts announcement-day price reactions only in the FG era (β = −1.01, p = 0.005, R² = 32.8\%, N = 392 fund-meeting observations {[}TODO: confirm N from price table{]}). Orthogonalizing tone to Target+Path is a no-op (orth R² = 0.4\%) --- tone is near-orthogonal by construction, which answers the expert's instruction and is itself informative.

\textbf{Reading.} In the ZLB/FG regime the \emph{level} tool is stuck; language does the work --- so words move prices. Outside FG, statements confirm decisions already priced in. {[}TODO: add CBI price regressions; event-study figure.{]}

\section{5. Portfolios: The Flow Null (H4)}\label{portfolios-the-flow-null-h4}

\textbf{Design.} flows\_\{c,m\} \textasciitilde{} MP\_shock\_m + CBI\_shock\_m (same-month window; also Post and Diff windows, all in \texttt{flow\_results\_rerun\_v8.json}).

\textbf{Result.} Across all 7 classes: no significant flow response to either shock in the same-month window except government bonds (CBI +0.759, p = 0.037, N = 157, R² = 3.4\%); adding tone (raw or orthogonalized) leaves flows unmoved (p = 0.64); the rebuilt LLM score likewise (p = 0.17).

\textbf{Interpretation.} This null is \emph{theoretical} confirmation, not absence of evidence: under Pillar 3's timing structure, monthly flows aggregate post-price-adjustment rebalancing; announcement-day information is already in prices. Government bonds are the exception that maps to §6's de-risking story --- the safety asset absorbs the rotation. {[}TODO: power analysis --- what effect size could we detect?{]}

\section{6. Risk-Taking: De-Risking on Good News (H5)}\label{risk-taking-de-risking-on-good-news-h5}

\textbf{Design.} Risk ladder (ordinal 1--7 across classes) × CBI, with sign asymmetry (risk\_x\_mp\_pos / risk\_x\_mp\_neg) and ZLB interaction (MP × ZLB).

\textbf{Result.} (i) Risk × CBI = −0.066 (p = 0.036, N = 1099): positive CBI news → flows out of high-risk classes \emph{relative to} low-risk --- de-risking on good news. (ii) MP × ZLB = −1.51 (p = 0.005): unconventional-era MP shocks push flows strongly out of risky classes. (iii) The naive ``bad news → safety'' story fails: risk\_x\_mp\_neg is economically nil (β ≈ 0.006, p ≈ 0.98). {[}TODO: coefficient table from H5\_sign\_asymmetry; connect to APT risk-premium mechanism in §2.3.{]}

\textbf{Why this matters.} Standard risk-taking-channel readings (Borio-Zhu, Adrian-Shin) predict risk-taking on \emph{accommodative} news. We find the opposite sign pattern for CBI: the central bank's \emph{better} news compresses risk premia and triggers rule-based rotation to safety --- a portfolio-balance reading, not a sentiment reading. {[}TODO: verify against Bauer-Swanson risk-premium evidence.{]}

\section{7. Measurement: Dictionary Failure and the LLM Diagnostic}\label{measurement-dictionary-failure-and-the-llm-diagnostic}

\textbf{The rebuild.} The original LLM measure proxied statements with meeting metadata --- circular (metadata embed outcomes, not language). Rebuilt: full-text stance scoring, 161/161, multi-vendor committee with majority sign rule; audited with a five-layer hallucination-defense pipeline (retrieval consistency, fact-check, CROWN anchor-independence 0/30, prompt-variant vote 100\%, self-consistency 70\%; cross-vendor pairwise r = 0.89--0.97). One one-off scoring error (Operation Twist 2011: router +85 vs four vendors dovish −60\ldots−90) caught by the cross-model layer and corrected --- the audit is part of the methodology. {[}TODO: appendix on the audit protocol.{]}

\textbf{Score compression.} LLM vs LM tone r = +0.351 overall; on the 11 divergence statements the dictionary output is nearly constant (std 1.10) while LLM separates directions (std ≈ 70). Divergence events: Operation Twist {[}removed after audit --- see note{]}, QE3, 2015 calendar guidance, 2017 run-off, 2020 pandemic QE (×4), 2022-06 acceleration --- each a paradigm-shift episode; enriched 1.4× in FG∪ZLB. The dictionary does not mis-read these; it \emph{fails to differentiate} them.

\textbf{Positioning.} The LLM measure adds no price or flow predictiveness (p = 0.17 flows; CBI correlation +0.027, as theory requires). Its contribution is diagnostic: tone's FG price effect is measured in precisely the regime where the dictionary is least reliable --- reported side by side, per the expert's measurement-first discipline.

\section{8. Robustness}\label{robustness}

\begin{enumerate}
\def\labelenumi{\arabic{enumi}.}
\tightlist
\item
  Orthogonalized tone (Bauer-Swanson style): no-op for flows (near-orthogonal by construction); reported anyway.
\item
  Window rules: Same/Post/Diff flow windows; gap ≤ 3 rule; cap 0.5. {[}TODO: table{]}
\item
  Sign-asymmetry and ZLB interactions: full table.
\item
  LLM audit: five-layer protocol + cross-vendor committee (Appendix A).
\item
  {[}TODO: placebo dates; alternative tone dictionaries; subsample stability 2006-2012 vs 2013-2026.{]}
\end{enumerate}

\section{9. Conclusion}\label{conclusion}

Communication moves prices when language is the policy tool (FG); portfolios move later and through rotation, not reaction; the rotation is triggered by \emph{good} news (de-risking on positive information surprises); and the dictionary-based text measures that dominate the literature silently lose resolution exactly when communication regime shifts --- a measurement failure the inversion-rule framework both predicts and repairs.

{[}TODO: limitations --- fund-flow aggregation granularity; statement-only scope (minutes/speeches out); LLM vendor dependence mitigated by committee rule.{]}

\section{Appendix A --- LLM Audit Protocol}\label{appendix-a-llm-audit-protocol}

Five layers (ported from the NeuroSync/federated-platform defense): retrieval consistency, vector fact-check, CROWN anchor-independence, multi-node vote (prompt variants + cross-vendor), self-consistency. Verdicts on the 30-statement audit sample: Verified 13 / LikelyTrue 7 / Uncertain 10 / LikelyFalse 0 / Hallucination 0. Thinking-model caveat (reasoning tokens vs max\_tokens). Files: \texttt{llm\_hallucination\_audit.py}, \texttt{llm\_5way\_analysis.py}, \texttt{paper/llm\_hallucination\_audit\_report.md}.

\section{Appendix B --- Variable Codebook (anchor)}\label{appendix-b-variable-codebook-anchor}

See \texttt{paper/codebook.docx}. Naming discipline: ``shock'' = innovations only; ``surprise'' only in GSS-construction context; tone is never ``sentiment''; ``information shock'' only in JK-discussion context (quoted).
